% Vietnamese translation for OI Public Library
% Source-PDF: IOI中国国家候选队论文/国家集训队2015论文集.pdf
\documentclass[11pt]{article}
\usepackage{oipl-vn}
\usepackage{array}

\title{Tập luận văn đội tuyển quốc gia Trung Quốc năm 2015}
\author{Huấn luyện viên: Yu Linyun}
\date{China Computer Federation, 2015}

\begin{document}
\maketitle

\origtitle{Tập luận văn ứng viên đội tuyển quốc gia Trung Quốc Olympic Tin học năm 2015}
\sourcepath{IOI China candidate team papers / National training team 2015 collection.pdf}

\renewcommand{\abstractname}{Tóm tắt}
\begin{abstract}
PDF nguồn là tập hợp luận văn đội tuyển quốc gia Trung Quốc năm 2015. Tệp được
tạo bằng pdfTeX nhưng lớp văn bản không ánh xạ đúng về Unicode: ngay cả mục
lục khi trích xuất bằng \texttt{pdftotext} cũng trở thành mojibake. Bản dịch
này chuyển ngữ phần đầu sách, mục lục đọc được từ các trang đã render, và hai
bài đầu tiên của tập sách, được OCR từ ảnh trang.
\end{abstract}

\section{Thông tin đầu sách}

Tập sách có tiêu đề \emph{Tập luận văn ứng viên đội tuyển quốc gia Trung Quốc
Olympic Tin học năm 2015}. Huấn luyện viên ghi trên bìa là Yu Linyun. Đơn vị
xuất bản là China Computer Federation.

\paragraph{Ghi chú trích xuất.}
Do lớp văn bản của PDF nguồn không dùng được, các bài đã dịch được đọc từ ảnh
render bằng \texttt{pdftoppm} và OCR bằng \texttt{tesseract -l chi\_sim+eng},
sau đó đối chiếu trực quan ở các trang biên. Hiện bản dịch đã hoàn thành hai
bài đầu tiên: bài ``Mở rộng suffix automaton trên trie'' dùng các trang PDF vật
lý 5--20, tương ứng trang in 1--16; bài ``Bàn về ứng dụng tư tưởng heuristic
trong thi tin học'' dùng các trang PDF vật lý 21--40, tương ứng trang in
17--36. Trang PDF vật lý 41, trang in 37, là đầu bài kế tiếp về khớp chuỗi.

\section{Mục lục đã dịch}

\begin{center}
\begin{tabular}{r >{\raggedright\arraybackslash}p{0.46\linewidth} >{\raggedright\arraybackslash}p{0.29\linewidth}}
\textbf{Trang} & \textbf{Bài viết} & \textbf{Tác giả / trường} \\
\hline
1 & Mở rộng suffix automaton trên trie & Liu Yanyi, Trường Yali, Trường Sa, Hồ Nam \\
17 & Bàn về ứng dụng tư tưởng heuristic trong thi tin học & Ren Zhizhou, Trường Trung học số 1 Thiệu Hưng \\
37 & Bàn về một số phương pháp khớp chuỗi & Wang Shenghao, Trường Trung học số 1 Thiệu Hưng \\
67 & Suffix automaton và ứng dụng & Zhang Tianyang, Trường Trung học số 1 Trường Sa, Hồ Nam \\
81 & Phép toán trên hàm sinh và bài toán đếm tổ hợp & Jin Ce, Trường Học Quân Hàng Châu, Chiết Giang \\
103 & Báo cáo ra đề \emph{ydc's bonus} & Liu Jiancheng, Trường Yali, Trường Sa \\
113 & Bàn về ứng dụng chia khối trong một lớp bài toán trực tuyến & Zou Daoyao, Trường Trung học Trấn Hải, Ninh Ba \\
129 & Thuật toán liên quan tới cactus và ứng dụng & Wang Yisong, Trường Học Quân Hàng Châu, Chiết Giang \\
149 & Bàn về thuật toán matching trên đồ thị và ứng dụng & Chen Haobo, Trường Trường Quận, Trường Sa, Hồ Nam \\
175 & Bàn về các bài toán vật lý trong thi tin học & Chen Siyu, Trường Trung học số 2 Hàng Châu \\
199 & Báo cáo ra đề \emph{Đề bài bị thất lạc} & Yu Jiping, Trường Trung học số 24 Đại Liên \\
205 & Một số kỹ thuật tối ưu DP & Zhang Hengjie, Trường Trung học số 1 Thiệu Hưng \\
217 & Báo cáo ra đề \emph{Product} & Du Yuhao, Trường Trung học Trấn Hải, Ninh Ba \\
231 & Bước đầu tìm hiểu một lớp bài toán lấy mã nguồn làm đầu vào & Lu Xiaochen, Học viện Giao Xuyên Trấn Hải, Ninh Ba \\
271 & Tính chất, ứng dụng và thuật toán nhanh cho lũy thừa tập hợp & Lu Kaifeng, Trường Trung học số 2 Vũ Hán, Hồ Bắc
\end{tabular}
\end{center}

\section{Mở rộng suffix automaton trên trie}

\begin{center}
\textbf{Liu Yanyi}\\
Trường Yali, Trường Sa, Hồ Nam
\end{center}

\subsection*{Tóm tắt}

Bài viết phân tích chi tiết một automaton hữu hạn xác định có thể nhận ra mọi
hậu tố của một trie. Tác giả đưa ra và chứng minh các kết luận rằng số trạng
thái là tuyến tính, tổng số hàm chuyển không vượt quá số nút của trie nhân với
kích thước bảng chữ cái, đồng thời trình bày một thuật toán xây dựng khả thi,
có độ phức tạp tuyến tính khi được phép xử lý offline. Bài viết cũng phân tích
độ phức tạp của thuật toán xây dựng online, giới thiệu một số ứng dụng liên
quan của suffix automaton, với hy vọng gợi mở thêm các hướng sử dụng cấu trúc
này.

\subsection{Lời nói đầu}

Trong thi thuật toán, xử lý xâu là một mảng rất quan trọng. Nhìn chung có hai
loại bài toán lớn: bài toán có mẫu cố định và bài toán có văn bản chính cố
định. Khi mẫu cố định, có thể có nhiều văn bản chính khác nhau. Nếu chỉ có một
mẫu, thuật toán KMP nổi tiếng \([2,3]\) xử lý được lớp bài toán này; cũng có thể
dùng hashing xâu \([13]\) để đếm số lần xuất hiện của mẫu. Nếu có nhiều mẫu,
thuật toán AC automaton quen thuộc \([1]\) đảm nhiệm tốt lớp bài toán đó. Nếu
văn bản chính cố định, ta thường xây suffix array \([11]\), suffix tree \([4,5]\)
hoặc suffix automaton \([7]\) cho văn bản chính. Lớp bài toán này đã được nghiên
cứu đầy đủ trong cộng đồng thi thuật toán. Chẳng hạn, Chen Lijie đã nêu nhiều
ứng dụng của suffix automaton trên một xâu trong bài giảng giao lưu trại đông
\([8]\), còn Luo Suijian giải thích chi tiết việc dùng suffix array trong thi
tin học trong luận văn đội tuyển quốc gia của mình \([12]\). Khi đề bài cho
phép xử lý offline, hai loại bài toán trên thường có thể chuyển đổi cho nhau,
thể hiện tính linh hoạt của các thuật toán xử lý xâu.

Trong kỳ chọn đội tuyển tỉnh Chiết Giang năm 2015 xuất hiện bài toán sau:
cho một cây có \(n\le 10^5\) nút, số lá không quá \(20\), mỗi cạnh mang một ký
tự. Đường đi giữa hai nút trên cây tạo thành một xâu; hỏi có tất cả bao nhiêu
xâu con khác nhau.\footnote{ZJOI 2015 Day 1, \emph{Vùng đất huyễn tưởng được
chư thần ưu ái}, tác giả Chen Lijie.} Bài này rõ ràng thuộc loại văn bản chính
cố định, nhưng văn bản chính không còn là một xâu mà là một cây. Nếu trực tiếp
liệt kê mọi cặp đỉnh rồi dùng trie \([1,13]\) để kiểm tra hai xâu có bằng nhau
hay không, độ phức tạp là \(O(n^3)\). Ngay cả dùng hashing xâu \([13]\), thời
gian vẫn là \(O(n^2)\), không thể qua được dữ liệu.

Ta chuyển hóa bài toán như sau. Vì chỉ có \(20\) lá, với mỗi lá làm gốc, dựng
một trie riêng rồi hợp nhất tất cả các trie này thành một trie mới. Khi đó mọi
xâu xuất hiện trên cây ban đầu chắc chắn xuất hiện trong trie mới; văn bản chính
đã trở thành một trie.

Để giải quyết vấn đề trên, bài viết phân tích automaton hữu hạn xác định nhận
ra toàn bộ hậu tố của một trie và đề xuất một cách xây dựng. Nội dung chính gồm
năm phần:
\begin{enumerate}
  \item đưa một số khái niệm trên xâu sang trie, đồng thời giới thiệu suffix
  automaton của trie;
  \item phân tích suffix link và xây dựng cấu trúc cây parent;
  \item chứng minh số trạng thái là tuyến tính và nêu chặn trên cho tổng số hàm
  chuyển trong automaton;
  \item mô tả thuật toán xây suffix automaton cho trie; khi được phép offline,
  độ phức tạp của nó là tuyến tính theo kích thước automaton;
  \item trình bày một số ứng dụng của suffix automaton trên trie, qua đó làm rõ
  hơn cách nhìn về suffix automaton.
\end{enumerate}

\subsection{Suffix automaton của trie}

Trong bài viết này, mọi ký tự của xâu và mọi ký tự trên cạnh của trie đều thuộc
một bảng chữ cái hữu hạn \(A\). Tập mọi xâu sinh bởi bảng chữ cái này là
\(A^*\); xâu rỗng được ký hiệu là \(\varepsilon\). Tập không chứa xâu rỗng là
\(A^+=A^*\setminus\{\varepsilon\}\). Dùng \(|s|\) để chỉ độ dài xâu \(s\),
\(s_i\) là ký tự thứ \(i\), và \(s_{i,j}\) là xâu gồm các ký tự từ vị trí \(i\)
đến vị trí \(j\) của \(s\). Ký hiệu \(a\cdot b\) là xâu thu được khi nối \(a\)
với \(b\); dấu nhân có thể được lược bỏ.

Tập mọi hậu tố của một xâu \(s\) là
\[
  S(s)=\{s_{i,|s|}\mid 1\le i\le |s|\}.
\]
Với một tập xâu \(S\), \(Minl(S)\) và \(Maxl(S)\) lần lượt là độ dài của xâu
ngắn nhất và dài nhất trong \(S\).

Dùng \(T\) để chỉ một trie, và dùng \(T_x,T_y,T_z,\ldots\) để chỉ các nút của
\(T\). Xâu nhận được bằng cách nối theo thứ tự các ký tự trên đường đi từ nút
\(T_x\) đến một nút \(T_y\) trong cây con của nó được ký hiệu là \(T_{x,y}\).
Nếu không nói rõ, khi dùng \(T_{x,y}\), ta luôn ngầm hiểu \(T_y\) nằm trong cây
con của \(T_x\).

Một tiền tố của \(T\) là
\[
  P_{T_x}=T_{\mathrm{root},x}\qquad
  (\mathrm{root}\text{ là gốc của }T,\ T_x\in T).
\]
Tập mọi xâu con của \(T\) là
\[
  F(T)=\{T_{x,y}\mid T_x,T_y\in T,\ T_y\text{ nằm trong cây con của }T_x\}.
\]
Tương tự, tập mọi hậu tố của \(T\) là
\[
  S(T)=\{T_{x,y}\mid T_x,T_y\in T,\ T_y\text{ là nút lá}\}.
\]

Xây suffix automaton cho trie nghĩa là xây một automaton hữu hạn có thể nhận ra
mọi xâu trong \(S(T)\), đồng thời cố gắng giảm số trạng thái và số chuyển xuống
càng nhỏ càng tốt. Tương tự suffix automaton trên một xâu, ta định nghĩa vị trí
xuất hiện của một xâu trên trie:
\[
  Right(s)=\{T_y\mid T_y\in T,\ \exists T_x\in T,\ T_y
  \text{ nằm trong cây con của }T_x,\ T_{x,y}=s\}.
\]
Định nghĩa quan hệ tương đương \(R_T\); khi không gây nhầm lẫn viết là \(R\):
\[
  aRb \Longleftrightarrow Right(a)=Right(b)\qquad (a,b\in A^*).
\]
Gọi \(R(s)\) là tập mọi xâu tương đương với \(s\), tức một lớp tương đương. Trong
automaton, mỗi trạng thái tương ứng với một lớp tương đương.

Gọi suffix automaton xây từ \(T\) là \(\Phi(T)\), trạng thái đầu là
\(\Phi(T)_\varepsilon\). Ký hiệu \(\Phi(T)_s\), với \(s\in F(T)\), là trạng thái
sau khi đọc \(s\). Hàm chuyển được định nghĩa bởi chuyển nhãn \(a\) từ
\(\Phi(T)_s\) tới \(\Phi(T)_{sa}\) khi \(a\in A\) và \(sa\in F(T)\). Khi nối hai
xâu \(s_1,s_2\) mà \(s_1,s_2,s_1s_2\in F(T)\), ta cũng có
\(\Phi(T)_{s_1}\) đọc \(s_2\) tới \(\Phi(T)_{s_1s_2}\). Theo định nghĩa,
\[
  \Phi(T)_a=\Phi(T)_b \Longleftrightarrow aRb \qquad (a,b\in F(T)).
\]
Nếu \(a\) là một xâu tương ứng với trạng thái \(\Phi(T)_b\), thì tập xâu tương
ứng với trạng thái \(\Phi(T)_b\) chính là \(R(b)\).

\begin{center}
\fbox{\begin{minipage}{0.82\linewidth}
\textbf{Hình a.} Trong trie minh họa của bản gốc, \(b\) tương đương với \(ab\);
các nút màu đỏ là tập \(Right\) của chúng.
\end{minipage}}
\end{center}

Để phân tích số trạng thái của toàn bộ cấu trúc, tức số lớp tương đương \(R\)
trong \(F(T)\), ta cần xét một số tính chất của các trạng thái suffix automaton.
Mục tiêu là làm số trạng thái nhỏ nhất có thể.

\paragraph{Bổ đề 2.1.}
Với \(x,y,z,u\in F(T)\), ta có:
\[
\begin{aligned}
&\text{(i)}\quad xRy \Rightarrow (x\in S(y)\text{ hoặc }y\in S(x)),\\
&\text{(ii)}\quad x\in S(z) \Rightarrow Right(z)\subseteq Right(x),\\
&\text{(iii)}\quad (x\in S(z),\ z\in S(u),\ xRu)\Rightarrow zRuRx.
\end{aligned}
\]

\paragraph{Chứng minh.}
(i) Không mất tính tổng quát, giả sử \(|x|\le |y|\) và
\(Right(x)=Right(y)\). Lấy tùy ý \(T_v\in Right(x)\). Khi đó
\(x\in S(P_{T_v})\) và \(y\in S(P_{T_v})\), nên \(x\in S(y)\).

(ii) Với mọi \(T_v\in Right(z)\), ta có \(z\in S(P_{T_v})\). Vì \(x\in S(z)\),
suy ra \(x\in S(P_{T_v})\), do đó \(T_v\in Right(x)\). Vậy
\(Right(z)\subseteq Right(x)\).

(iii) Từ (ii), \(z\in S(u)\) cho \(Right(u)\subseteq Right(z)\), và
\(x\in S(z)\) cho \(Right(z)\subseteq Right(x)\). Lại có \(xRu\), nên
\(Right(z)=Right(u)=Right(x)\), tức \(zRuRx\).

Mệnh đề (i) cũng cho thấy trong cùng một lớp tương đương \(R\), không thể có hai
xâu khác nhau nhưng cùng độ dài.

\paragraph{Hệ quả 2.2.}
Với một lớp tương đương \(R(a)\), gọi \(x\) là xâu dài nhất trong \(R(a)\), và
\(y\) là xâu ngắn nhất trong \(R(a)\). Khi đó
\[
  R(a)=\{z\mid z\in S(x)\text{ và } |y|\le |z|\le |x|\}.
\]
Điều này suy ra trực tiếp từ Bổ đề 2.1.

\subsection{Suffix link}

Trong thuật toán xây suffix automaton cho một xâu, suffix link giữ vai trò then
chốt \([7,6]\), và toàn bộ thuật toán rất giống KMP \([3]\). Khi xây suffix
automaton cho trie, ta cũng giữ suffix link và so sánh quá trình xây dựng với
thuật toán xây automaton tối tiểu xác định nhận ra tập \(A^*\{P_{T_x}\}\), tức
cách xây AC automaton quen thuộc \([1]\). Tương tự hàm \(\mathrm{fail}\) trong
AC automaton, ta định nghĩa suffix link của suffix automaton trên trie như sau.

Suffix link là một hàm \(s_T:A^*\setminus\{\varepsilon\}\to A^*\), khi không
nhầm lẫn viết là \(s\):
\[
  s_T(x)=\{z\mid z\in S(x),\ |z|=Minl(R_T(x))-1\}.
\]
Nói cách khác, đó là hậu tố dài nhất \(y\) của \(x\) sao cho \(xR_Ty\) không
đúng.

\paragraph{Bổ đề 3.1.}
Với \(x,y\in F(T)\), ta có:
\[
\begin{aligned}
&\text{(i)}\quad xRy\Rightarrow s(x)=s(y),\\
&\text{(ii)}\quad Right(x)\subset Right(s(x)),\\
&\text{(iii)}\quad Minl(R(x))-1=|s(x)|=Maxl(R(s(x))).
\end{aligned}
\]

\paragraph{Chứng minh.}
(i) Không mất tính tổng quát, giả sử \(|x|\le |y|\). Theo định nghĩa,
\(|s(x)|=Minl(R(x))-1\), \(|s(y)|=Minl(R(y))-1\), mà \(R(x)=R(y)\), nên
\(|s(x)|=|s(y)|<|x|\le |y|\). Theo Bổ đề 2.1(i), \(x\in S(y)\). Lại có
\(s(x)\in S(x)\) và \(s(y)\in S(y)\), do đó \(s(x)=s(y)\).

(ii) Theo Bổ đề 2.1(ii), \(Right(x)\subseteq Right(s(x))\). Vì \(xRs(x)\) là sai,
hai tập này không bằng nhau, nên bao hàm là nghiêm ngặt.

(iii) Vì \(s(x)\in R(s(x))\), ta có \(|s(x)|\le Maxl(R(s(x)))\). Giả sử tồn tại
xâu \(y\) tương đương với \(s(x)\) và dài nhất trong lớp đó. Khi ấy
\(Right(y)=Right(s(x))\), đồng thời \(s(x)\in S(y)\). Dùng phản chứng: nếu
\(|y|>|s(x)|\), lấy \(T_v\in Right(x)\). Từ (ii), \(T_v\in Right(s(x))\), suy ra
\(T_v\in Right(y)\), nên \(y\in S(P_{T_v})\) và \(x\in S(P_{T_v})\). Nếu
\(|y|>|x|\) thì \(x\in S(y)\), theo Bổ đề 2.1(ii) sẽ có
\(Right(s(x))=Right(y)\subseteq Right(x)\), mâu thuẫn với (ii). Nếu
\(|y|\le |x|\), vì \(xRy\) không đúng nên \(s(x)=y\), mâu thuẫn với
\(|y|>|s(x)|\). Vậy mệnh đề được chứng minh.

Từ (i), mọi xâu trong cùng một lớp tương đương \(R(a)\) có suffix link trỏ tới
cùng một xâu. Vì vậy ta có thể gán suffix link cho từng trạng thái của suffix
automaton. Cụ thể, suffix link của \(R(a)\) trỏ tới \(R(s(a))\). Theo (ii) hoặc
(iii), các suffix link này tạo thành một cây; ta gọi nó là \term{cây parent}.
Định nghĩa \(Fa(\Phi(T)_s)\) là suffix link của trạng thái \(\Phi(T)_s\), tức
\[
  Fa(\Phi(T)_a)=\Phi(T)_b \Longleftrightarrow s(a)Rb.
\]

\paragraph{Bổ đề 3.2.}
Với \(x\in F(T)\), lấy tùy ý \(T_v\in Right(x)\). Nếu
\(|P_{T_v}|>Maxl(R(x))\) và đặt \(a=P_{T_v}\), thì
\[
  s\bigl(a_{|a|-Maxl(R(x)),\,|a|}\bigr)Rx.
\]

\paragraph{Chứng minh.}
Đặt \(b=a_{|a|-Maxl(R(x)),\,|a|}\). Từ giả thiết có \(x\in S(b)\). Theo Hệ quả
2.2, vì \(|b|>Maxl(R(x))\), nên \(bRx\) không đúng, còn hậu tố độ dài
\(Maxl(R(x))\) của \(b\) lại tương đương với \(x\). Vì vậy \(b\) không tương
đương với hậu tố đó; theo định nghĩa suffix link, \(s(b)\) chính là hậu tố này,
và \(s(b)Rx\).

\begin{center}
\fbox{\begin{minipage}{0.82\linewidth}
\textbf{Hình b.} Cây trong bản gốc là cây parent của suffix automaton được xây
từ trie ở Hình a.
\end{minipage}}
\end{center}

\subsection{Số trạng thái tuyến tính}

Với suffix automaton xây cho một xâu \(x\), kích thước automaton là \(O(|x|)\)
\([7,6]\). Ta cũng mong suffix automaton xây cho trie có quy mô trạng thái
tuyến tính. Sau khi định nghĩa suffix link, cần phân tích kỹ hơn \(\Phi(T)\).

\paragraph{Bổ đề 4.1.}
Với \(x,y\in F(T)\), ta có:
\[
\begin{aligned}
&\text{(i)}\quad \forall y,\ \neg(xRs(y)) \Rightarrow |Right(x)|=1,\\
&\text{(ii)}\quad |Right(x)|=1 \Rightarrow P_{T_v}Rx
   \quad (Right(x)=\{T_v\}),\\
&\text{(iii)}\quad \forall T_v\in Right(x),\ \neg(xRP_{T_v})\\
&\hspace{5em}\Rightarrow
  \exists u,v,\ s(u)Rs(v)Rx\text{ và }\neg(uRv).
\end{aligned}
\]

\paragraph{Chứng minh.}
(i) Chứng minh bằng phản chứng. Giả sử \(|Right(x)|>1\). Lấy
\(T_{v_1},T_{v_2}\in Right(x)\) và đặt \(a=P_{T_{v_1}}, b=P_{T_{v_2}}\). Khi đó
\(|a|,|b|\ge Maxl(R(x))\); nếu một trong hai độ dài nhỏ hơn \(Maxl(R(x))\), xâu
trong \(R(x)\) có độ dài \(Maxl(R(x))\) không thể xuất hiện ở nút tương ứng, mâu
thuẫn. Nếu \(a\notin R(x)\), theo Bổ đề 3.2 suy ra suffix link của một hậu tố
thích hợp của \(a\) trỏ vào lớp của \(x\). Nếu \(a\in R(x)\), vì
\(T_{v_1}\ne T_{v_2}\), ta có \(a\ne b\), và \(|b|\ge Maxl(R(x))\), nên một lập
luận tương tự cũng cho một xâu có suffix link trỏ vào \(R(x)\). Mệnh đề được
chứng minh.

(ii) Biết \(x\in S(P_{T_v})\), do đó \(|Right(P_{T_v})|\le |Right(x)|\). Vì
\(|Right(x)|=1\), suy ra \(Right(P_{T_v})=\{T_v\}=Right(x)\), tức \(P_{T_v}Rx\).

(iii) Từ phản đảo của (i), nếu \(|Right(x)|\ne 1\) thì \(|Right(x)|\ge 2\).
Theo giả thiết, với mọi \(T_v\in Right(x)\) ta có
\(|P_{T_v}|>Maxl(R(x))\). Tồn tại \(T_{v_1},T_{v_2}\) sao cho hậu tố chung dài
nhất của \(P_{T_{v_1}}\) và \(P_{T_{v_2}}\) có độ dài đúng bằng \(Maxl(R(x))\).
Nếu không, mọi cặp đều có hậu tố chung dài hơn \(Maxl(R(x))\); theo Bổ đề
2.1(iii), hậu tố chung dài hơn đó cũng tương đương với \(x\), mâu thuẫn với
định nghĩa \(Maxl(R(x))\).

Đặt
\[
  a=(P_{T_{v_1}})_{|P_{T_{v_1}}|-Maxl(R(x)),\,|P_{T_{v_1}}|},\qquad
  b=(P_{T_{v_2}})_{|P_{T_{v_2}}|-Maxl(R(x)),\,|P_{T_{v_2}}|}.
\]
Từ cách chọn, \(a\ne b\), \(a\notin S(b)\), \(b\notin S(a)\). Theo Bổ đề 2.1(i),
\(\neg(aRb)\). Dùng Bổ đề 3.2 ta có \(s(a)Rx\) và \(s(b)Rx\). Lấy \(u=a\),
\(v=b\) là được.

\paragraph{Định lý 4.2.}
Suffix automaton xây cho một trie có số trạng thái tuyến tính.

\paragraph{Chứng minh.}
Xét cây parent tạo bởi suffix link. Từ Bổ đề 4.1(i) và (ii), với mọi lá
\(\Phi(T)_s\) trên cây parent đều tồn tại một tiền tố \(P_{T_v}\) sao cho
\(P_{T_v}Rs\). Từ Bổ đề 4.1(iii), mọi trạng thái \(\Phi(T)_s\) mà với mọi
\(T_v\) đều không có \(sRP_{T_v}\) thì phải có ít nhất hai con trong cây parent.
Nói cách khác, các lá và các nút chỉ có một con trên cây parent đều tương ứng
với một tiền tố nào đó \(P_{T_v}\); các nút còn lại có ít nhất hai con.

Số tiền tố chỉ là \(|T|\), và trạng thái đạt được khi đọc từng \(P_{T_v}\) là
duy nhất. Vì vậy tổng số lá và nút một con trên cây parent không vượt quá
\(|T|\); toàn bộ số nút của cây không vượt quá \(2|T|\). Do đó số trạng thái là
tuyến tính.

Nếu không chỉ số trạng thái mà tổng số hàm chuyển, tức số cạnh, cũng tuyến tính,
thì toàn bộ automaton sẽ tuyến tính. Tuy nhiên, khác với suffix automaton của
một xâu, tổng số hàm chuyển của suffix automaton trên trie có thể ở cấp số trạng
thái nhân với kích thước bảng chữ cái. Ta chứng minh nó không vượt quá cấp này,
và đưa ra một trie đạt cùng cấp độ.

\paragraph{Định lý 4.3.}
Với một trie, chặn trên của số hàm chuyển trong suffix automaton là
\[
  O((\text{số nút của trie})\cdot(\text{kích thước bảng chữ cái})).
\]

\paragraph{Chứng minh.}
Gọi trie là \(T\). Theo Định lý 4.2, \(|\Phi(T)|\le 2|T|\). Với mỗi trạng thái
\(\Phi(T)_s\), số chuyển nhiều nhất là \(|A|\). Vì vậy tổng số hàm chuyển không
vượt quá \(2|T|\,|A|\).

Ta cũng dựng được một trie khiến số hàm chuyển đạt chặn này theo cấp độ. Chỉ cần
xét những ký tự thực sự xuất hiện trên cạnh của \(T\), nên \(|A|\le |T|\). Một
phần của \(T\) là một dây dài \(|T|/2\), mọi cạnh trên dây đều mang cùng ký tự
\(a\). Các nút còn lại nối từ nút cuối của dây bằng các ký tự đôi một khác
nhau; gọi tập ký tự này là \(B\). Khi đó \(a^i b\in S(T)\) với
\(0\le i\le |T|/2\) và \(b\in B\), đồng thời hiển nhiên \(\neg(a^iRa^j)\) nếu
\(i\ne j\). Vì thế với mỗi trạng thái \(\Phi(T)_{a^i}\), mọi \(b\in B\) đều tạo
một hàm chuyển. Tổng số hàm chuyển đạt \(|T|/2\cdot |B|\), tức cùng cấp với
\(|T|/2\cdot |A|/2\), trùng cấp với chặn trên trong Định lý 4.3.

\begin{center}
\fbox{\begin{minipage}{0.82\linewidth}
\textbf{Hình c.} Ví dụ trie trong bản gốc gồm một dây toàn ký tự \(a\), rồi tỏa
ra nhiều cạnh ký tự khác nhau, làm tổng số hàm chuyển đạt chặn trên theo cấp độ.
\end{minipage}}
\end{center}

Dù số chuyển khi xây suffix automaton cho trie không luôn tuyến tính, trong phần
lớn đề thi và bài toán thực tế, bảng chữ cái không lớn; thường chỉ là chữ cái
Latin thường hoặc chữ số. Khi bảng chữ cái là hằng số, số chuyển vẫn tuyến tính,
nên đây vẫn là một phương pháp hiệu quả để xử lý các bài toán xâu liên quan đến
trie.

\subsection{Thuật toán xây dựng}

Trong thuật toán xây suffix automaton cho xâu \([7,8]\), ta dùng một thuật toán
tăng dần online: từ suffix automaton của \(s\), xây suffix automaton của
\(sa\), với \(a\in A\). Khi xây cho trie, ta cũng xét thuật toán tăng dần
online. Từ phân tích số trạng thái và số chuyển ở trên, ta biết cận dưới độ
phức tạp của thuật toán xây dựng là \(|T|\,|A|\).

Giả sử từ trie \(T\), ta thêm vào nút \(T_v\) một cạnh nhãn \(a\) đi tới nút con
mới \(T_w\), thu được trie mới \(T_N\). Đặt \(s=P_{T_v}\), nên tiền tố mới là
\(sa=P_{T_w}\). Xét những phần của \(\Phi(T_N)\) cần sửa trên nền \(\Phi(T)\).
Trước hết phân tích số trạng thái.

\subsubsection{Thay đổi trạng thái}

Với mọi xâu đã có thể tiếp tục mở rộng thành hậu tố của trie mới, trạng thái
tương ứng trong \(\Phi(T_N)\) phải tồn tại. Ngược lại, nếu một xâu không thể mở
rộng thành hậu tố mới, trạng thái đó không cần xuất hiện thêm. Nếu một lớp tương
đương không đổi sau khi thêm \(T_w\), trạng thái suffix automaton tương ứng cũng
không đổi, tức \(\Phi(T_N)_x=\Phi(T)_x\).

\paragraph{Trường hợp 1.}
Nếu \(T_v\) đã có một con nối bằng cạnh nhãn \(a\), ta không cần sửa
\(\Phi(T)\).

\paragraph{Trường hợp 2.}
Nếu \(sa\in F(T)\), nghĩa là trong \(\Phi(T)\) đã có chuyển nhãn \(a\) từ trạng
thái của \(s\), thì lớp tương đương \(R_T(sa)\) có thể bị tách làm hai lớp.

Thật vậy, với \(x,y\in F(T)\) và \(xR_Ty\), nếu trong tập \(Right_{T_N}\) của
\(x\) và \(y\) đều xuất hiện nút mới \(T_w\), thì khi bỏ \(T_w\) ra, hai tập
vẫn bằng nhau, nên chắc chắn \(xR_{T_N}y\). Vì vậy mọi xâu tương đương với
\(sa\) trong \(T_N\) vốn đều nằm trong cùng lớp \(R_T(sa)\). Với
\(x\in R_T(sa)\), nếu \(xR_{T_N}sa\), không có trạng thái mới. Nếu tồn tại
\(x\in R_T(sa)\) nhưng \(\neg(xR_{T_N}sa)\), cần tách \(R_T(sa)\) thành hai lớp:
một lớp tương đương với \(sa\) trong \(T_N\), một lớp không tương đương với nó.
Lớp không tương đương có tập \(Right\) không chứa \(T_w\), nên tập \(Right\) của
nó giữ nguyên như trong \(T\).

Xét suffix link của hai trạng thái mới. Gọi \(u,v\in F(T)\), trong đó
\(uR_{T_N}sa\), \(\neg(vR_{T_N}sa)\), nhưng \(uR_Tsa\) và \(vR_Tsa\). Nếu đặt
\(u\) vào lớp thứ nhất và \(v\) vào lớp thứ hai, vì \(sa=P_{T_w}\) và
\(uR_{T_N}sa\), mọi hậu tố của \(u\) nhận thêm cùng vị trí mới nên suffix link
của lớp thứ nhất không đổi. Với \(v\), phần mới bị tách ra đúng tại lớp của
\(sa\), do đó suffix link của lớp thứ hai trỏ tới lớp mới chứa \(sa\). Nói ngắn
gọn, nếu một trạng thái bị tách, một bản sao giữ suffix link cũ, bản sao còn
lại có suffix link trỏ tới trạng thái của \(sa\).

\begin{center}
\fbox{\begin{minipage}{0.82\linewidth}
\textbf{Hình d.} Suffix automaton của trie trong Hình a.

\textbf{Hình e, f, g.} Sau khi thêm một nút vào trie của Hình a, bản gốc minh
họa lần lượt trie mới, suffix automaton mới và cây parent mới. Nút và cạnh màu
cam biểu thị phần mới trong automaton/cây parent.
\end{minipage}}
\end{center}

Khi đã biết lớp tương đương chứa \(sa\) trong \(T_N\), mọi hậu tố của \(sa\)
trước đó đã được đưa vào các lớp liên quan, và các trạng thái ấy không đổi, nên
không cần xét lại.

\paragraph{Trường hợp 3.}
Nếu \(sa\notin F(T)\), ta tạo một lớp mới \(R_{T_N}(sa)\). Vì xuất hiện một hậu
tố mới, vẫn có khả năng một lớp tương đương cũ bị tách.

Có một kết luận khá rõ ràng: \(s_{T_N}(sa)\) bằng hậu tố dài nhất của \(sa\) có
xuất hiện trong \(F(T)\). Chứng minh đơn giản: nếu \(x\in S(sa)\) nhưng
\(x\notin F(T)\), thì \(x\) chỉ xuất hiện tại nút mới \(T_w\), và \(sa\) cũng
chỉ xuất hiện tại \(T_w\); do đó \(Right(sa)=Right(x)=\{T_w\}\), nên \(xRsa\).
Ngược lại, nếu \(x\in S(sa)\) và \(x\in F(T)\), thì \(|Right(x)|>1\), nên
\(\neg(xRsa)\).

Kết luận này giúp quá trình xây dựng thuận tiện hơn. Các hậu tố
\[
  y\in S(sa),\qquad y\notin S(s_{T_N}(sa))
\]
tạo thành lớp mới \(R_{T_N}(sa)\), và ta tạo một trạng thái mới cho lớp đó.
Với những \(y\in S(s_{T_N}(sa))\), tức phần hậu tố còn lại, có thể xem
\(s_{T_N}(sa)\) như \(sa\) trong Trường hợp 2. Dù Trường hợp 2 dùng tính chất
\(sa=P_{T_w}\), bản thân \(s_{T_N}(sa)\) đã có tính cực đại cần thiết, nên không
gây sai lệch.

Bài toán còn lại là tìm \(s_{T_N}(sa)\). Theo kết luận trên, ta cần hậu tố dài
nhất của \(sa\) đã xuất hiện ít nhất một lần trong \(F(T)\). Bắt đầu từ trạng
thái \(\Phi(T)_s\), đi lên cây parent qua các suffix link cho đến khi gặp trạng
thái đầu tiên có chuyển nhãn \(a\). Nếu trạng thái đó đại diện bởi một xâu
\(r\), thì \(ra\) là hậu tố dài nhất của \(sa\) xuất hiện ít nhất hai lần, tức
\(s_{T_N}(sa)=ra\).

Từ ba trường hợp trên, khi thêm một ký tự vào trie, số trạng thái tăng nhiều
nhất hai. Điều này một lần nữa chứng minh số trạng thái tuyến tính, đồng thời
mô tả được thay đổi của suffix link. Ta vẫn chưa xét cách sửa hàm chuyển; phần
tiếp theo xử lý điểm này.

\subsubsection{Thay đổi hàm chuyển}

Trong Trường hợp 2, nếu không tạo trạng thái mới, chắc chắn không có thay đổi
hàm chuyển. Nếu có, lớp cũ \(R_T(sa)\) bị tách thành hai lớp \(R_{T_N}(sa)\) và
một lớp còn lại. Vì nút mới \(T_w\) là lá, nó không làm thay đổi các chuyển vốn
có từ các xâu cũ. Các chuyển của hai trạng thái mới có thể sao chép từ trạng
thái cũ tương ứng.

Tuy nhiên, vì lớp \(R_T(sa)\) cũ bị xóa, có thể tồn tại các trạng thái có chuyển
nhãn \(a\) trỏ tới trạng thái cũ đó; những chuyển này phải được sửa. Theo phân
tích ở Trường hợp 2, các chuyển dọc theo chuỗi suffix link từ trạng thái của
\(s\) lên trên, chừng nào chúng còn đi tới lớp cũ, sẽ được đổi để trỏ tới lớp
mới phù hợp. Những chuyển khác giữ nguyên. Cách làm trực tiếp là giữ một bản
sao đại diện cho phần không đổi, tạo một trạng thái mới cho phần bị tách, rồi
bạo lực sửa các chuyển nhãn \(a\) trên chuỗi suffix link cho tới khi gặp trạng
thái mà chuyển nhãn \(a\) không còn trỏ tới lớp cũ.

Trong Trường hợp 3, ta tạo trạng thái mới \(R_{T_N}(sa)\); phần còn lại giống
Trường hợp 2. Vì \(T_w\) là lá, trạng thái mới ban đầu không có chuyển ra. Do có
trạng thái mới, một số trạng thái trên chuỗi suffix link có thể cần thêm chuyển
nhãn \(a\) trỏ tới trạng thái mới của \(sa\). Theo phân tích Trường hợp 3, chỉ
cần thêm các chuyển này dọc theo chuỗi từ trạng thái của \(s\) đi lên cho tới
khi gặp trạng thái đã có chuyển nhãn \(a\) thích hợp.

\subsubsection{Phân tích độ phức tạp}

Tương tự trường hợp một xâu, toàn bộ thuật toán là thuật toán tăng dần online.
Với trie \(T\), số lần thực hiện là \(|T|\). Mỗi lần chỉ thêm nhiều nhất hai
trạng thái, nên phần trạng thái là \(O(|T|)\). Suffix link của mỗi trạng thái
mới chỉ sửa một lần, phần này cũng là \(O(|T|)\).

Xét hàm chuyển. Trong Trường hợp 2 có thể cần sao chép một tập chuyển; mỗi trạng
thái có nhiều nhất \(|A|\) chuyển, nên phần này là \(O(|T|\,|A|)\). Phần phức
tạp hơn là các lần bạo lực sửa chuyển khác. Có thể chứng minh tổng thời gian
cho phần này không vượt quá tổng độ sâu các lá trong trie; gọi tổng đó là
\(G(T)\).

Ý tưởng chứng minh giống phân tích thế năng. Gọi \(Dep(y)\), với \(y\in F(T)\),
là độ sâu của trạng thái \(\Phi(T)_y\) trong cây parent. Xét đường từ trạng thái
\(\Phi(T)_s\) và trạng thái mới tới gốc trên cây parent. Theo Bổ đề 2.1(iii),
nếu tồn tại lớp của \(sa\) thì cũng tồn tại lớp của \(s\), nên
\(Dep(sa)\le Dep(s)+1\). Khi bạo lực sửa chuyển, ta liên tục đặt
\(c=s(c)\) và sửa chuyển của \(c\); mỗi bước như vậy làm \(Dep(c)\) giảm \(1\).
Khi dừng, độ sâu của trạng thái liên quan chỉ có thể tăng thêm một hằng số so
với \(Dep(s)\). Vì mỗi lần sửa đều làm giảm thế năng, tổng số lần sửa là cấp độ
độ sâu.

Có thể xem \(Dep\) là hàm thế năng của trạng thái hiện tại. Với mỗi lá của
\(T\), nếu thực hiện riêng quá trình chèn theo đường từ gốc tới lá đó, mỗi phép
chèn làm \(Dep\) tăng nhiều nhất \(3\), còn mỗi lần bạo lực sửa chuyển làm
\(Dep\) giảm \(1\). Do đó số lần bạo lực sửa chuyển thực chất bị chặn bởi độ
sâu; cộng trên mọi lá sẽ được kết luận trên.

Độ phức tạp này chưa phải đẹp, nhưng độ phức tạp của xây dựng online có thể đạt
đến cấp đó.

\begin{center}
\fbox{\begin{minipage}{0.82\linewidth}
\textbf{Hình h.} Nếu chèn các nút theo thứ tự chỉ số nút trong ví dụ của bản
gốc, độ phức tạp đạt \(O(G(T))\).
\end{minipage}}
\end{center}

Khi không yêu cầu online, ta có thể thêm các nút trên trie theo thứ tự FIFO, tức
thứ tự duyệt rộng. Khi đó chứng minh độ phức tạp của J. A. Blumer và cộng sự
\([9,10]\) cho thuật toán xây trên xâu cũng áp dụng cho trie. Chỉ những tổ hợp
xâu con cố định mới khiến ta phải bạo lực sửa hàm chuyển \([6]\), gợi ý rằng
phần độ phức tạp này không vượt quá \(O(|T|)\).

Tóm lại, nếu yêu cầu online, độ phức tạp xây dựng là \(O(G(T))\). Nếu cho phép
offline, ta có thể xây theo thứ tự FIFO và đạt độ phức tạp
\(O(|T|\,|A|)\), đúng với cận dưới đã nêu.

\subsection{Ứng dụng}

\subsubsection{ZJOI 2015: Vùng đất huyễn tưởng được chư thần ưu ái}

Đề bài đã xuất hiện ở phần lời nói đầu, nên không nhắc lại nữa.

Sau chuyển hóa, bài toán trở thành: đếm số xâu con khác nhau về bản chất của
một trie. Bằng phương pháp xây suffix automaton cho trie ở trên, ta thu được
một automaton có thể nhận ra mọi hậu tố của trie. Mỗi xâu con của trie chắc
chắn là tiền tố của một hậu tố nào đó, giống hệt tính chất khi xây suffix
automaton cho một xâu. Vì vậy, với một trạng thái \(\Phi(T)_s\), số xâu con khác
nhau mà nó đại diện là \(|R(s)|\); cộng \(|R(s)|\) trên mọi trạng thái sẽ cho
đáp án.

Cũng có cách khác. Sau khi xây automaton, trực tiếp đếm có bao nhiêu xâu khác
nhau không bị automaton từ chối; điều này tương đương đếm số xâu con khác nhau,
hay đếm số đường chuyển khác nhau tồn tại khi bắt đầu từ trạng thái đầu. Xem
automaton như một đồ thị có hướng không chu trình, ta chỉ cần quy hoạch động đơn
giản trên đồ thị này.

Theo phân tích độ phức tạp ở trên, đề bài không yêu cầu online, nên có thể dùng
thuật toán offline với độ phức tạp \(O(n|A|)\). Thực ra điều kiện số lá là hằng
số rất mạnh; có thể chứng minh automaton xây bằng phương pháp trên có thời gian
\(O(n)\). Lý do là ta có thể tách riêng đường từ gốc đến từng lá thành các xâu,
rồi dùng phân tích độ phức tạp giống trường hợp xâu \([9]\). Vì điều này không
liên quan trực tiếp đến nội dung chính của bài viết, tác giả bỏ qua chứng minh
chi tiết.

Đến đây bài toán trong phần lời nói đầu đã được giải quyết trọn vẹn.

Từ bài toán này có thể thấy hầu hết bài toán xâu xử lý được bằng suffix
automaton đều có thể mở rộng lên trie: tìm xâu con thứ \(K\) nhỏ nhất của trie,
tìm xâu con chung dài nhất của nhiều trie, đếm số lần xuất hiện của xâu vòng
trên trie, v.v. Cũng có những bài rất dễ trên xâu nhưng phức tạp hơn khi mở
rộng lên trie, chẳng hạn đếm số xâu khác nhau mà tiền tố của chúng là một xâu
cho trước. Các bài này đều có thể giải sau khi xây suffix automaton cho trie.

\subsubsection{So sánh với AC automaton}

Trong các thuật toán xâu quen thuộc, AC automaton xây một trie từ \(m\) mẫu rồi
xây các con trỏ fail, sau đó có thể nhận diện số lần xuất hiện của từng mẫu
trong xâu đầu vào. Ta có thể thay thế cách dùng này bằng suffix automaton trên
trie: xây trie từ \(m\) xâu, rồi xây suffix automaton cho trie đó. Khi đọc xâu
cần nhận diện, nếu trạng thái hiện tại \(\Phi(T)_s\) không có chuyển nhãn \(a\),
ta nhảy tới \(Fa(\Phi(T)_s)\) và tiếp tục nhận diện. Rõ ràng tổng số lần nhảy
không vượt quá độ dài xâu đang nhận diện.

Về độ phức tạp, trie \(T\) xây từ \(m\) xâu có \(G(T)\) không vượt quá tổng độ
dài của \(m\) xâu, nên độ phức tạp không kém hơn AC automaton. Nhưng thuật toán
xây ở đây có thể là online, trong khi thuật toán xây AC automaton chỉ là
offline.

\subsubsection{Cây parent}

\begin{center}
\fbox{\begin{minipage}{0.82\linewidth}
\textbf{Hình i.} Bản gốc minh họa việc bổ sung các ký tự trên cạnh của Hình b
để nhìn cây parent như một trie của các tiền tố đảo.
\end{minipage}}
\end{center}

Thực chất, ta có thể đảo toàn bộ tiền tố trong trie rồi xây lại thành một trie,
tương tự suffix tree của một xâu. Với mỗi trạng thái \(\Phi(T)_s\), sắp xếp
các con của nó trong cây parent. Cụ thể, nếu
\[
  Fa(\Phi(T)_x)=\Phi(T)_s,
\]
ta sắp xếp các trạng thái \(\Phi(T)_x\) theo thứ tự của đoạn
\[
  x_{1,\,|x|-Maxl(R(s))}.
\]
Khi đó ta thu được một mảng tiền tố của trie, tức thứ tự sau khi đảo mọi tiền tố.
Ta cần duy trì cho \(\Phi(T)_s\) một nút \(T_v\in Right(s)\); kết hợp nhân đôi
trên cây, có thể nhanh chóng tìm được đoạn \(x_{1,\,|x|-Maxl(R(s))}\).

Với tính chất này, có thể nhanh chóng xác định trạng thái tương ứng khi thêm một
ký tự \(a\in A\) vào trước xâu đang nhận diện. Giả sử hiện ở trạng thái
\(\Phi(T)_s\). Ta tìm nhị phân trong các con của nó trên cây parent để tìm
\(\Phi(T)_{as}\), hoặc kết luận \(\Phi(T)_{as}=\Phi(T)_s\).

\subsection{Lời cảm ơn}

Xin cảm ơn thầy Qu Yunhua của Trường Yali đã lâu dài hướng dẫn và quan tâm tới
việc học tập cũng như cuộc sống của tôi.

Xin cảm ơn cha mẹ đã nuôi dưỡng tôi.

Xin cảm ơn các thầy Zhu Quanmin và Wang Xingming của Trường Yali đã chỉ dạy.

Xin cảm ơn CCF đã cung cấp nền tảng và cơ hội giao lưu.

Xin cảm ơn các huấn luyện viên đội tuyển quốc gia vì sự tận tụy.

Xin cảm ơn các anh Huang Zhixiang của Đại học Thanh Hoa và Huang Yuyang của Đại
học Giao thông Thượng Hải đã giúp đỡ tôi.

Xin cảm ơn các bạn Yang Dingcheng, Kuang Zhengfei và Liu Jiancheng của Trường
Yali đã đọc duyệt bài viết này.

Xin cảm ơn tất cả những bạn học từng giúp đỡ tôi.

\subsection*{Tài liệu tham khảo của bài viết}

\begin{enumerate}[label={[\arabic*]},leftmargin=3em]
  \item A. V. Aho and M. J. Corasick, Efficient string matching: An aid to bibliographic research.
  \item R.-S. Boyer and J. S. Moore, A fast string searching algorithm.
  \item D. E. Knuth, J. H. Morris and V. I. Pratt, Fast pattern-matching in strings.
  \item P. Weiner, Linear pattern-matching algorithms.
  \item E. M. McCreight, A space-economical suffix-tree construction algorithm.
  \item M. Mohri, P. Moreno and E. Weinstein, General Suffix Automaton Construction Algorithm and Space Bounds.
  \item M. Crochemore, Transducers and repetitions.
  \item Chen Lijie, \emph{Suffix automaton}.
  \item J. A. Blumer, Algorithms for the directed acyclic word graph and related structures.
  \item A. Blumer, J. Blumer, D. Haussler, R. M. McConnell and A. Ehrenfeucht, Complete inverted files for efficient text retrieval and analysis.
  \item U. Manber and G. Myers, Suffix arrays: A new method for on-line string searches.
  \item Luo Suijian, \emph{Suffix array: một công cụ mạnh để xử lý xâu}.
  \item Thomas H. Cormen, Charles E. Leiserson, Ronald L. Rivest and Clifford Stein, \emph{Introduction to Algorithms}.
\end{enumerate}

\section{Bàn về ứng dụng tư tưởng heuristic trong thi tin học}

\begin{center}
\textbf{Ren Zhizhou}\\
Trường Trung học số 1 Thiệu Hưng
\end{center}

\subsection*{Tóm tắt}

Thuật toán heuristic là các thuật toán được thiết kế dựa trên kinh nghiệm. So
với các thuật toán tối ưu hóa truyền thống, heuristic thường đánh đổi độ chính
xác hoặc tính ổn định để lấy tốc độ. Bài viết thử đưa ra một định nghĩa cụ thể
hơn cho heuristic trong thi tin học, đồng thời thảo luận một số thuật toán
heuristic thường dùng và một số mô hình bài toán. Trong đó có một thuật toán
xấp xỉ cho cây Steiner nhỏ nhất với độ phức tạp \(O(|S||V|^2)\); bài viết cũng
chứng minh tỉ lệ xấp xỉ của thuật toán này đạt
\[
  2\left(1-\frac1{\mathrm{leaf}}\right)
  \le 2\left(1-\frac1{|S|}\right),
\]
trong đó \(\mathrm{leaf}\) là số lá của cây Steiner tối ưu.

\subsection{Dẫn nhập}

Thuật toán heuristic là một lớp thuật toán được nêu ra trong tương quan với các
thuật toán tối ưu hóa. Chúng được xây dựng từ trực giác hoặc kinh nghiệm, với
mục tiêu tìm được một nghiệm khá tốt trong chi phí thời gian và bộ nhớ chấp
nhận được. Khái niệm này rất rộng; đa số heuristic khá phức tạp và không thể
đưa thẳng vào bài thi OI.

Tác giả giản lược và cải tiến một số thuật toán heuristic để chúng phù hợp hơn
với OI, rồi định nghĩa thêm hai kiểu heuristic trong bối cảnh thi tin học:
\begin{itemize}
  \item dựa vào trực giác để điều chỉnh và tối ưu thuật toán, hy sinh tính ổn
  định để đổi lấy tốc độ;
  \item dựa vào trực giác để trực tiếp dựng nghiệm tốt, hy sinh tính đúng tuyệt
  đối để giải nhanh.
\end{itemize}

Nhiều thuật toán quen thuộc cũng chứa tư tưởng heuristic. Sau tối ưu, chúng
không làm giảm độ phức tạp xấu nhất, nhưng hiệu suất thực tế thường tăng rõ rệt,
chẳng hạn:
\begin{itemize}
  \item truy vấn điểm gần nhất trên \texttt{k-d tree};
  \item Bellman--Ford tối ưu bằng hàng đợi, tức SPFA, và các tối ưu tham lam
  liên quan;
  \item mô phỏng tôi luyện, thuật toán di truyền và các thuật toán tìm kiếm
  ngẫu nhiên lấy cảm hứng từ khoa học tự nhiên.
\end{itemize}

Bài viết thảo luận ba mô hình heuristic:
\begin{itemize}
  \item thuật toán \(A^*\) và các biến thể, dùng để nối giữa thuật toán xấp xỉ
  và lời giải tối ưu;
  \item phân loại Bayes ngây thơ, một thuật toán xác suất kết hợp phán đoán
  trực quan với suy luận toán học;
  \item gộp heuristic, một tối ưu đơn giản cho các bài toán cần hợp nhất dữ
  liệu.
\end{itemize}

Ngoài việc áp dụng mô hình kinh điển, nhiều khi ta cũng có thể tự thiết kế
heuristic. Bài viết dành nhiều dung lượng cho hai bài toán cổ điển:
\begin{itemize}
  \item bài toán chia đều có trọng số, hay \emph{2-partition problem}, từ đó có
  một thuật toán xấp xỉ có xác suất đúng rất cao và quy trình thiết kế đáng tham
  khảo;
  \item bài toán cây Steiner, được trình bày kỹ hơn, với một thuật toán hiệu quả
  có thể đánh giá sai số xấp xỉ, rồi kết hợp với \(A^*\) để thu được thuật toán
  xấp xỉ cho \(K\) cây Steiner nhỏ nhất.
\end{itemize}

Ở mỗi phần, tác giả đều đưa ví dụ với hy vọng giúp người đọc nắm được cách sử
dụng và tiếp tục phát triển các ý tưởng này.

\subsection{\texorpdfstring{Thuật toán \(A^*\)}{Thuật toán A*}}

Trong một số bài toán rất khó tìm thuật toán đa thức hiệu quả, ta buộc phải
dùng tìm kiếm, trong khi số trạng thái lại rất lớn. Khi đó có thể thiết kế hàm
ước lượng để cắt tỉa; phần này trình bày một số cách dùng ước lượng để tăng tốc
tìm kiếm.

\subsubsection{Tìm kiếm heuristic}

Giả sử trạng thái hiện tại là \(s\), chi phí thực tế đã đi tới \(s\) là
\(g(s)\), ước lượng chi phí từ \(s\) tới trạng thái đích là \(h(s)\), còn chi
phí tối ưu thật sự từ \(s\) tới đích là \(h^*(s)\). Đặt
\[
  F(s)=g(s)+h(s).
\]
Ta có thể dựa vào giá trị này, mỗi lần mở rộng trạng thái có \(F(s)\) tốt nhất,
và cắt tỉa theo tính chất của hàm ước lượng. Trong bài viết, các bài toán được
mặc định là tối thiểu hóa chi phí.

\subsubsection{\texorpdfstring{\(h(s)=h^*(s)\)}{h(s)=h*(s)}}

Khi hàm ước lượng cho đúng chi phí còn lại, mỗi lần mở rộng trạng thái tốt nhất
theo \(F(s)\) sẽ nhanh chóng thu được nghiệm tối ưu.

Sau khi đã có nghiệm tối ưu, nếu tiếp tục mở rộng các trạng thái còn lại thì có
thể nhận được nghiệm tốt thứ \(K\), đồng thời tránh tính các trạng thái thừa.
Phần tổng quát và các tối ưu cho bài toán đặc biệt đã được Yu Yili trình bày
trong luận văn đội tuyển quốc gia năm 2014.

\subsubsection{\texorpdfstring{\(h(s)\le h^*(s)\)}{h(s)<=h*(s)}}

Đôi khi, để tính ước lượng thuận tiện, ta bỏ qua một số ràng buộc trong đề bài.
Kết quả thu được là một nghiệm thô ``tốt hơn'' nghiệm tối ưu thật, tức ước lượng
không vượt quá chi phí còn lại. Khi đó ước lượng vẫn cung cấp căn cứ cắt tỉa
đúng, giúp giảm phạm vi trạng thái phải xét.

\paragraph{IDA*.}
IDA*, hay tìm kiếm sâu dần lặp có heuristic, là một biến thể của \(A^*\) kết hợp
ý tưởng lặp sâu dần với hàm ước lượng. Trong các bài có số trạng thái lớn và khó
lưu trữ, cách này tiết kiệm rất nhiều bộ nhớ, đồng thời mã cài đặt cũng gọn.

Với lặp sâu dần đơn giản, mỗi lượt thường chọn một bước tăng cố định và dùng
giới hạn độ sâu để cắt tỉa. Với IDA*, ta cắt tỉa bằng ước lượng. Vì
\(h(s)\le h^*(s)\), giá trị nhỏ nhất của \(g(s)+h(s)\) trong các trạng thái bị
cắt có thể dùng để ước lượng đáp án kế tiếp, giảm số lần lặp so với bước tăng
cố định.

\paragraph{Ví dụ 1: UVa 10181, \emph{15-Puzzle Problem}.}
Bài toán 15-puzzle mở rộng từ 8-puzzle. Trên bảng \(4\times4\) có 15 ô đánh số
và một ô trống; mỗi bước có thể đưa một ô kề ô trống vào vị trí trống. Cần tìm
số bước ít nhất để tới trạng thái đích. Dữ liệu bảo đảm mọi trường hợp có lời
giải đều tìm được trong không quá 45 bước.

Việc phán đoán vô nghiệm cần phân tích riêng về puzzle số, bài viết không đi
sâu. Với phần tìm kiếm, một ước lượng rất đơn giản là tổng khoảng cách Manhattan
từ mỗi ô không trống tới vị trí đích. Cách này tương đương với việc bỏ qua ràng
buộc ``chỉ được đổi chỗ với ô trống'', đồng thời bỏ qua ảnh hưởng của một phép
đổi chỗ tới ô còn lại.

Vì đã bỏ bớt ràng buộc, ta có \(h(s)\le h^*(s)\). Với dữ liệu có lời giải trong
45 bước, hiệu quả cắt tỉa và hiệu quả giảm số vòng lặp đều khá tốt.

\subsubsection{Anytime algorithm}

Lớp này còn được gọi là thuật toán có thể ngắt giữa chừng. Không phải bài nào
cũng dễ tìm hàm ước lượng hiệu quả. Đôi khi ta chỉ ước lượng rất thô, không kiểm
soát được \(h(s)\le h^*(s)\), nên cũng không thể thu hẹp không gian tìm kiếm một
cách đúng đắn. Những bài như vậy thường xuất hiện dưới dạng nộp đáp án. Ta có
thể dùng thuật toán cải thiện dần, trong thời gian thi cố gắng thu được nghiệm
càng tốt càng tốt.

\paragraph{Ví dụ 2: CTSC 2009 Day 2, trò chơi số \(N\times N\).}
Trò chơi gồm \(N^2-1\) ô trượt và một ô trống. Mục tiêu là di chuyển ô trống
lên, xuống, trái, phải để đưa các ô về một thứ tự quy định. Đây là bài nộp đáp
án, với \(3\le N\le 6\) và 10 bộ dữ liệu.

So với ví dụ 1, miền dữ liệu lớn hơn nhiều. Đổi lại, bài không yêu cầu nghiệm
chính xác; chỉ cần trong thời gian cho phép tìm được nghiệm đủ tốt thì vẫn có
điểm đáng kể.

Nếu dùng trực tiếp ước lượng Manhattan của ví dụ 1 rồi chạy IDA*, ta có thể
nhanh chóng tìm được nghiệm tối ưu cho 5 bộ nhỏ. Nhưng giá trị ước lượng vẫn
cách đáp án thật khá xa, nên các bộ còn lại không xử lý được.

Một cách sửa đơn giản là đặt
\[
  h'(s)=c\,h(s),
\]
với \(c\) là hằng số điều chỉnh theo từng bộ dữ liệu. So với \(h(s)\),
\(h'(s)\) gần số bước thật hơn, nên có thể đưa vào \(A^*\). Tuy vậy, \(A^*\)
tốn bộ nhớ rất lớn, đó cũng là lý do ví dụ 1 không dùng \(A^*\).

Với bài nộp đáp án, ta có thể thỏa hiệp: chỉ giữ lại một phần trạng thái có ước
lượng tốt nhất. Mỗi lần chọn trạng thái tốt nhất để mở rộng; nếu không còn chỗ
lưu trạng thái mới thì bỏ trạng thái kém nhất. Vì bài có rất nhiều phương án
lời giải, việc loại bỏ một số trạng thái ``xấu'' ít ảnh hưởng đến đáp án cuối.
Ngoại trừ một bộ đặc biệt, cách này có thể đạt khoảng 9 đến 10 điểm mỗi bộ.

\subsection{Gộp heuristic}

Gộp heuristic là một chiến lược hợp nhất đơn giản, có thể tối ưu hiệu quả cho
các thuật toán cần gộp dữ liệu.

\subsubsection{Gộp heuristic trên cấu trúc dữ liệu}

Giả sử cần gộp hai cấu trúc dữ liệu \(A,B\), số phần tử đang được chúng quản lý
lần lượt là \(a,b\) và \(a\le b\). Ta tách toàn bộ thông tin trong cấu trúc nhỏ
\(A\), rồi lần lượt chèn chúng vào cấu trúc lớn \(B\).

\paragraph{Định lý 3.1.}
Giả sử sau một số phép gộp ta thu được một cấu trúc dữ liệu có \(n\) phần tử.
Khi dùng chiến lược gộp heuristic, tổng số lần chèn phát sinh là
\(O(n\log n)\).

\paragraph{Chứng minh.}
Xét riêng từng phần tử. Khi một phần tử bị chèn sang cấu trúc khác, kích thước
cấu trúc chứa nó sau phép gộp ít nhất gấp đôi trước đó. Sau \(O(\log n)\) lần
chèn, phần tử đó chắc chắn đã nằm trong cấu trúc cuối cùng. Cộng trên mọi phần
tử, ta được \(O(n\log n)\).

Chiến lược này có thể gộp các cấu trúc dữ liệu phức tạp, và cũng dễ mở rộng cho
các kiểu tập dữ liệu khác.

\paragraph{Ví dụ 3: HNOI 2009, \emph{Bánh pudding huyền ảo}.}
Có \(n\) chiếc pudding xếp thành một hàng và \(m\) thao tác. Mỗi lần đổi toàn bộ
pudding của một màu thành màu khác, rồi hỏi hiện có bao nhiêu đoạn màu liên
tiếp. Dữ liệu có \(n\le 10^6\).

Ta dùng danh sách liên kết hoặc \texttt{vector} để lưu các vị trí của từng màu.
Khi đổi màu, chọn danh sách nhỏ hơn để tách thô, sửa từng vị trí và cập nhật
đóng góp vào số đoạn. Theo Định lý 3.1, tổng số lần sửa là \(O(n\log n)\), nên
có thể bảo trì trực tiếp.

\paragraph{Ví dụ 4: BZOJ 3510, \emph{Thủ đô}.}
Cần duy trì một rừng \(n\) đỉnh với ba thao tác:
\begin{itemize}
  \item nối hai đỉnh \(u,v\), bảo đảm sau khi nối vẫn là rừng;
  \item hỏi trọng tâm của cây chứa đỉnh \(x\), nếu có hai trọng tâm thì lấy chỉ
  số nhỏ hơn;
  \item hỏi xor chỉ số trọng tâm của tất cả các cây.
\end{itemize}

Trước hết xét việc thêm một lá vào một cây. Trọng tâm chỉ có thể di chuyển về
phía lá mới, và nhiều nhất di chuyển một lần.

Khi nối hai cây, ta tách thô cây có ít đỉnh hơn và thêm từng đỉnh vào cây lớn
dưới dạng lá, đồng thời duy trì trọng tâm. Theo Định lý 3.1, số lần thêm lá là
\(O(n\log n)\); chỉ cần thêm một cấu trúc dữ liệu để kiểm tra trọng tâm có di
chuyển hay không.

\subsubsection{Hợp nhất theo hạng trong DSU}

Hợp nhất theo hạng là một cách duy trì DSU rất thực dụng: mỗi lần nối gốc của
tập nhỏ vào gốc của tập lớn. Từ Định lý 3.1 suy ra chiều cao cây là
\(O(\log n)\). Khi cài đặt, cũng có thể trực tiếp duy trì độ cao rồi hợp nhất
theo hạng.

Cách này bảo đảm độ sâu cây DSU ngay cả khi không nén đường đi. Trong một số
bài, nó tránh được việc mất thông tin khi nén đường đi.

\paragraph{Ví dụ 5: NOIP 2013, \emph{Vận tải bằng xe tải}.}
Có \(n\) thành phố, \(m\) đường hai chiều; mỗi đường có giới hạn trọng lượng
\(z_i\). Với \(q\) xe tải, cần biết mỗi xe có thể chở tối đa bao nhiêu mà không
vượt quá giới hạn trên đường đi. Dữ liệu:
\[
  n\le 10000,\quad m\le 50000,\quad q\le 30000,\quad z_i\le 100000.
\]

Cách thông thường là dùng Kruskal dựng cây khung lớn nhất. Giá trị nhỏ nhất trên
đường đi trong cây khung lớn nhất chính là giới hạn lớn nhất có thể giữa hai
đỉnh; có thể trả lời bằng nhân đôi, với độ phức tạp
\(O(m\log m+q\log n)\).

Nếu dùng hợp nhất theo hạng thay cho nén đường đi, ta có thể ghi lại với mỗi
đỉnh trong DSU trọng số cạnh nối nó với cha. Giá trị nhỏ nhất trên đường trong
cây DSU vẫn biểu diễn giới hạn rộng nhất giữa hai đỉnh. Trả lời thô có độ phức
tạp \(O(m\log m+q\log n)\), nhưng cài đặt ngắn hơn cách dựng cây rồi nhân đôi.

Vì bài toán chuyển hóa vẫn là truy vấn giá trị nhỏ nhất trên đường đi, ta vẫn có
thể dùng nhân đôi; về lý thuyết phần hỏi đạt \(O(\log\log n)\). Nếu chọn thuật
toán sắp xếp tuyến tính và vẫn nén đường đi trong DSU theo cách nối cạnh ban
đầu, độ phức tạp lý thuyết có thể đạt
\[
  O(n+m\alpha(n)+q\log\log n).
\]

\subsection{Bài toán trừu tượng và học máy}

Khi gặp các bài toán phân loại trừu tượng, cảm giác trực quan thường không cho
căn cứ phán đoán chính xác. Khi đó có thể dùng học máy thay cho quan sát thủ
công. Phần này thông qua một bài Google Code Jam để giới thiệu ngắn gọn một
thuật toán phân loại thực dụng.

\subsubsection{Bayes ngây thơ}

\paragraph{Công thức Bayes.}
Ký hiệu \(P(A\mid B)\) là xác suất sự kiện \(A\) xảy ra khi đã biết \(B\) xảy
ra; ký hiệu \(P(A\cap B)\) là xác suất cả \(A\) và \(B\) cùng xảy ra.

\paragraph{Định lý 4.1.}
\[
  P(A\mid B)=\frac{P(B\mid A)P(A)}{P(B)}.
\]

\paragraph{Chứng minh.}
Ta có
\[
  P(A\cap B)=P(A)P(B\mid A)=P(B)P(A\mid B).
\]
Chuyển vế sẽ thu được công thức Bayes.

\paragraph{Phán đoán ngây thơ.}
Gọi \(S\) là sự kiện cần phân loại, \(A\) và \(B\) là hai lớp có thể chọn. Nếu
\(P(A\mid S)>P(B\mid S)\), ta cho rằng \(S\) thuộc lớp \(A\); ngược lại thuộc
lớp \(B\). Khi có nhiều lớp hơn, chọn lớp có xác suất lớn nhất.

\paragraph{Giả định ngây thơ.}
Giả sử sự kiện cần phân loại có các thuộc tính
\[
  S=\{S_1,S_2,\ldots,S_m\}.
\]
Nếu giả định các thuộc tính đặc trưng này độc lập với nhau khi đã biết lớp, ta
có thể xấp xỉ
\[
  P(S\mid A)\approx \prod_{i=1}^m P(S_i\mid A).
\]
So sánh các xác suất xấp xỉ như vậy sẽ cho kết quả phân loại. Các giá trị
\(P(S_i\mid A)\) có thể tính bằng lấy mẫu ngẫu nhiên hoặc bằng một thuật toán
riêng cho bài toán.

\subsubsection{Ví dụ 6: GCJ Round 1A 2014, \emph{Proper Shuffle}}

Có hai cách sinh hoán vị. Đề cho 120 hoán vị độ dài 1000, mỗi hoán vị được sinh
độc lập bởi một trong hai thuật toán. Cần nhận diện hoán vị nào do thuật toán
nào sinh ra; nếu trả lời đúng ít nhất 109 trong 120 hoán vị thì được chấp nhận.

\begin{verbatim}
Algorithm 1 GOOD
for k = 0 to n - 1 do
  a[k] = k
end for
for k = 0 to n - 1 do
  p = random_int(k..n - 1)
  swap(a[k], a[p])
end for

Algorithm 2 BAD
for k = 0 to n - 1 do
  a[k] = k
end for
for k = 0 to n - 1 do
  p = random_int(0..n - 1)
  swap(a[k], a[p])
end for
\end{verbatim}

Thuật toán \texttt{GOOD} sinh mọi hoán vị với xác suất bằng nhau, còn
\texttt{BAD} thì không.

\paragraph{Phân tích đơn giản.}
Với \texttt{GOOD}, sau khi quyết định vị trí \(k\), xác suất xuất hiện của mọi
giá trị tại vị trí đó đều là \(1/n\), và các bước sau không ảnh hưởng vị trí đó.

Với \texttt{BAD}, khi quyết định phép đổi tại vị trí \(k\), nếu
\texttt{swap(a[k], a[p])} có \(p<k\), giá trị vốn ở vị trí \(a[k]\) bị đổi tới
một vị trí cao hơn, rồi sau đó còn có thể tiếp tục bị đổi tới vị trí cao hơn
nữa. Điều này tạo ảnh hưởng tinh vi tới xác suất một giá trị xuất hiện tại từng
vị trí.

\paragraph{Một đặc trưng đơn giản.}
Gọi \(S\) là hoán vị đầu vào. Có thể dùng số lượng vị trí thỏa \(S_i<i\) làm đặc
trưng. Phân tích của cuộc thi cho thấy phân bố đặc trưng này giữa \texttt{GOOD}
và \texttt{BAD} khác nhau rõ rệt; nếu dùng một ngưỡng thích hợp thì tỉ lệ đúng
có thể vượt 90\%.

\paragraph{Phân loại Bayes ngây thơ.}
Ta muốn tính \(P(\texttt{GOOD}\mid S)\). Theo Định lý 4.1:
\[
\begin{aligned}
P(\texttt{GOOD}\mid S)
&=\frac{P(S\mid\texttt{GOOD})P(\texttt{GOOD})}
{P(S\mid\texttt{GOOD})P(\texttt{GOOD})+
P(S\mid\texttt{BAD})P(\texttt{BAD})}.
\end{aligned}
\]
Vì \(P(\texttt{GOOD})=P(\texttt{BAD})\), công thức rút gọn thành
\[
  P(\texttt{GOOD}\mid S)=
  \frac{P(S\mid\texttt{GOOD})}
  {P(S\mid\texttt{GOOD})+P(S\mid\texttt{BAD})},
\]
và tương tự
\[
  P(\texttt{BAD}\mid S)=
  \frac{P(S\mid\texttt{BAD})}
  {P(S\mid\texttt{GOOD})+P(S\mid\texttt{BAD})}.
\]
Nếu \(P(\texttt{GOOD}\mid S)>P(\texttt{BAD}\mid S)\), xác suất \(S\) do
\texttt{GOOD} sinh ra lớn hơn. Điều kiện này tương đương
\[
  P(S\mid\texttt{GOOD})>P(S\mid\texttt{BAD}).
\]

Với \texttt{GOOD}, mọi hoán vị cùng xác suất:
\[
  P(S\mid\texttt{GOOD})=\frac1{n!}.
\]
Việc tính chính xác \(P(S\mid\texttt{BAD})\) lại rất khó, nên ta xấp xỉ:
\[
  P(S\mid\texttt{BAD})\approx
  \prod_{k=0}^{n-1}P(S_k\mid\texttt{BAD}).
\]
Để so sánh công bằng hơn, ta cũng dùng xấp xỉ tương tự cho \texttt{GOOD}:
\[
  P(S\mid\texttt{GOOD})\approx \prod_{k=0}^{n-1}\frac1n.
\]

Nếu \(P_{i,j}\) là xác suất sau khi chạy \texttt{BAD} ta có \(S_i=j\), ta dễ
thiết kế một DP \(O(n^2)\) để tính toàn bộ bảng này. Từ đó có thể tính
\(P(S\mid\texttt{BAD})\) xấp xỉ rồi so sánh với
\(P(S\mid\texttt{GOOD})\). Khi tính trực tiếp xác suất có thể gặp yêu cầu độ
chính xác cao; thực tế có thể so sánh \(P(S\mid\texttt{BAD})\,n^n\) với \(1.0\),
hoặc dùng ngưỡng khác vì bản thân quá trình đã là xấp xỉ.

\subsection{Bài toán chia đều có trọng số}

\paragraph{Ví dụ 7.}
Cho một đa tập số nguyên dương \(S\). Ký hiệu \(f(S)\) là tổng các phần tử trong
\(S\). Chia \(S\) thành hai tập \(S_1,S_2\), với \(S_1+S_2=S\), sao cho
\(|f(S_1)-f(S_2)|\) nhỏ nhất, đồng thời đưa ra phương án chia.

\subsubsection{Ví dụ đơn giản}

Với \(S=\{3,1,1,2,2,1\}\), có hai cách chia hợp lệ làm
\(f(S_1)=f(S_2)=5\), tức \(|f(S_1)-f(S_2)|=0\):
\[
  S_1=\{1,1,1,2\},\quad S_2=\{2,3\};
\]
\[
  S_1=\{3,1,1\},\quad S_2=\{2,2,1\}.
\]

\subsubsection{Một thuật toán quy hoạch động}

Có thể thiết kế DP giống ba lô: đặt \(d[i][j]\) cho biết sau khi dùng \(i\) phần
tử đầu tiên, có thể đạt \(f(S_1)=j\) hay không. Độ phức tạp là
\(O(nf(S))\).

Khi mọi phần tử trong \(S\) đều nhỏ, đây đúng là thuật toán tốt. Nhưng thực chất
đây chỉ là thuật toán giả đa thức.

\subsubsection{Một ý tưởng tham lam đơn giản}

Ý tưởng tự nhiên là đọc lần lượt từng phần tử và đặt nó vào tập hiện có tổng nhỏ
hơn. Để hai tổng dần áp sát nhau, ta sắp xếp các phần tử giảm dần trước khi xử
lý; độ phức tạp là \(O(n\log n)\).

Thuật toán này rất tự nhiên và kết quả thường không tệ, nhưng phân tích kỹ sẽ
thấy nó có lỗ hổng rõ ràng.

\subsubsection{Thuật toán sai phân}

Thuật toán sai phân cải tiến thuật toán tham lam ở trên. Ý tưởng ngắn gọn,
nhưng chất lượng nghiệm xấp xỉ tốt hơn nhiều.

\paragraph{Tư tưởng cơ bản.}
Khi đặt hai số \(X,Y\) vào hai tập khác nhau, phần ảnh hưởng tới đáp án chỉ là
độ lớn của hiệu giữa chúng. Điều này tương đương tạo ra một số mới
\(|X-Y|\), rồi tiếp tục xem nó như một số thường trong thuật toán.

Nhìn lại thuật toán tham lam trước đó, xét trạng thái ở mỗi quyết định. Giả sử
phần tử mới là \(X\), còn các phần tử đã quyết định tạo ra sai phân \(Y\). Việc
đưa \(X\) vào tập có tổng nhỏ hơn tương đương thay \(X,Y\) bằng
\(|X-Y|\). Nhưng mục đích ban đầu của ta là xử lý phần tử lớn trước, nên quyết
định theo thứ tự đọc vào là chưa tối ưu.

\paragraph{Quy trình thuật toán.}
Theo tư tưởng sai phân, ta dùng heap và DSU để bảo trì:
\begin{enumerate}
  \item Đưa mọi phần tử của \(S\) vào heap. Trong DSU, mỗi phần tử ban đầu là một
  tập riêng.
  \item Lấy hai phần tử lớn nhất khỏi heap, gọi là \(X,Y\); trong DSU chúng tương
  ứng với hai tập \(S_X,S_Y\). Quyết định đặt \(S_X\) và \(S_Y\) ở hai phía khác
  nhau, đưa \(|X-Y|\) trở lại heap, và nối hai tập trong DSU bằng một cạnh trọng
  số \(1\) biểu diễn việc đổi phía.
  \item Nếu heap chỉ còn một phần tử thì dừng, ngược lại lặp bước trước.
  \item Chọn gốc của cây DSU thuộc \(S_1\). Các phần tử khác dựa vào tính chẵn lẻ
  của độ dài đường tới gốc để quyết định thuộc \(S_1\) hay \(S_2\): chẵn cùng
  phía, lẻ khác phía.
\end{enumerate}

Khi cài đặt DSU, có thể duy trì tính chẵn lẻ của độ dài đường trong lúc nén
đường đi, hoặc dùng hợp nhất theo hạng để bảo đảm chiều cao \(O(\log n)\).

\paragraph{Hiệu quả xấp xỉ.}
Thuật toán có độ phức tạp \(O(n\log n)\), rất hiệu quả; điều quan trọng là nghiệm
có đủ tốt hay không.

Khi \(n\) không quá nhỏ, chẳng hạn \(n>1000\), với tập số sinh ngẫu nhiên
\(S_i\in[0,2^{30})\), thuật toán gần như chắc chắn chia được \(S_1,S_2\) sao cho
\[
  |f(S_1)-f(S_2)|\le 1.
\]
Nói cách khác, nếu \(f(S)\) lẻ thì sai khác bằng \(1\), còn nếu \(f(S)\) chẵn
thì sai khác bằng \(0\). Có thể thấy xác suất đúng của thuật toán này rất cao.

Khi \(n\) nhỏ, có thể dùng thuật toán này để thiết kế hàm ước lượng rồi chạy tìm
kiếm heuristic nhằm thu được nghiệm tốt hơn.

\subsubsection{Ví dụ 8: POI 2013, \emph{Polarization}}

Cho một cây \(n\) đỉnh. Hãy định hướng mỗi cạnh thành cạnh có hướng. Nếu từ đỉnh
\(u\) đi được tới đỉnh \(v\), gọi \((u,v)\) là một cặp hợp lệ. Cần tìm số cặp
hợp lệ nhỏ nhất và lớn nhất trong mọi cách định hướng. Dữ liệu có
\(n\le 250000\).

\paragraph{Chuyển hóa bài toán.}
Cây là đồ thị hai phía, nên giá trị nhỏ nhất chắc chắn là \(n-1\).

Với giá trị lớn nhất, nghiệm tối ưu luôn tồn tại một đỉnh \(x\) sao cho mọi đỉnh
còn lại hoặc có đường đi tới \(x\), hoặc có đường đi từ \(x\) tới nó. Phân tích
tiếp cho thấy \(x\) phải chọn ở trọng tâm của cây.

Lấy \(x\) làm gốc, cây tách thành một số cây con. Bài toán còn lại là chia các
cây con này thành hai tập sao cho tích kích thước của hai tập lớn nhất.

\paragraph{Thuật toán tối ưu truyền thống.}
Chú ý rằng tổng các kích thước ở đây là \(n\), nên số giá trị phần tử khác nhau
chỉ là \(O(\sqrt n)\). Ta có thể chuyển thành bài toán ba lô đa trọng, một bài
toán kinh điển không trình bày lại ở đây. Độ phức tạp là \(O(n\sqrt n)\).

\paragraph{Thuật toán xấp xỉ.}
Nếu dùng trực tiếp thuật toán sai phân, có thể qua khoảng 80\% dữ liệu. Nguyên
nhân chính là khi số phần tử ít, sai lệch của nghiệm xấp xỉ lớn hơn.

Ta đặt một ngưỡng. Khi số cây con nhỏ hơn ngưỡng này, dùng ba lô vét cạn; các
trường hợp còn lại dùng thuật toán sai phân. Cách ghép này có thể qua toàn bộ
dữ liệu chính thức của POI, với độ phức tạp \(O(n\log n)\).

Khi thuật toán tối ưu truyền thống đủ nhanh, không nên ưu tiên thuật toán xấp
xỉ. Ví dụ này chỉ nhằm cho thấy thuật toán sai phân có thể tạo ra nghiệm xấp xỉ
rất tốt.

\subsection{Bài toán cây Steiner nhỏ nhất}

\paragraph{Ví dụ 9.}
Cho đồ thị vô hướng có trọng số \(G=(V,E)\). Tập đỉnh đặc biệt \(S\) là một tập
con của \(V\). Cần tìm một đồ thị liên thông
\[
  T=(V_T,E_T)
\]
sao cho \(S\subseteq V_T\subseteq V\), \(E_T\subseteq E\),
\(|E_T|=|V_T|-1\), và tổng trọng số các cạnh trong \(E_T\) là nhỏ nhất.

\subsubsection{Một thuật toán tối ưu}

Cây Steiner nhỏ nhất có thuật toán trạng thái nén kinh điển, chuyển về mô hình
đường đi ngắn nhất; có thể xem luận văn đội tuyển quốc gia năm 2014 của Jiang
Jinye. Độ phức tạp là
\[
  O(|V|\,3^{|S|}+|E|\,2^{|S|}),
\]
đây là thuật toán tốt nhưng chỉ áp dụng được cho miền dữ liệu nhỏ.

\subsubsection{Một ý tưởng tham lam đơn giản}

Một tham lam đơn giản là tìm cây khung nhỏ nhất của đồ thị gốc, rồi xóa các cạnh
thừa trên cây khung, tức xóa nhiều nhất có thể nhưng vẫn giữ cho tập \(S\) liên
thông.

Thuật toán này có lỗ hổng lớn. Trong ví dụ của bản gốc, với
\[
  V=\{A,B,C,1,2,3,4,5\},\quad S=\{A,B,C\},\quad x>2,
\]
tham lam theo cây khung nhỏ nhất chọn các cạnh trọng số \(x\) và cho nghiệm
kém, trong khi nghiệm tối ưu đi qua các cạnh trọng số \(1\) và các cạnh
\(x+1\) như hình minh họa.

\subsubsection{Thuật toán heuristic}

Dựa trên ví dụ trên, ta cải tiến có chủ đích.

\paragraph{Phân tích thêm.}
Trong ví dụ, nếu tạm bỏ qua các điểm 1 và 2, rồi xét khoảng cách ngắn nhất giữa
các điểm còn lại, cây khung nhỏ nhất trên đồ thị rút gọn đã có thể cho nghiệm
tối ưu. Để tổng quát hơn, ta loại cả những điểm không thuộc \(S\) khỏi đồ thị
rút gọn.

Vấn đề là các đường đi ngắn nhất giữa các cặp đỉnh của \(S\) có thể giao nhau.
Khi cuối cùng tạo cây, các cạnh giao nhau không được tính lặp.

\paragraph{Quy trình thuật toán.}
Sau khi chỉnh lý, ta có thuật toán sau:
\begin{enumerate}
  \item Từ đồ thị gốc \(G=(V,E)\), dựng đồ thị đầy đủ
  \(G_1=(S,E_1)\). Trọng số cạnh giữa hai đỉnh trong \(S\) là độ dài đường đi
  ngắn nhất giữa chúng trong \(G\).
  \item Tìm cây khung nhỏ nhất \(T_1\) của \(G_1\). Nếu có nhiều nghiệm thì chọn
  tùy ý.
  \item Dựng một đồ thị con \(G_2=(V_2,E_2)\) của \(G\): với mỗi cạnh của \(T_1\),
  đưa vào \(G_2\) một đường đi ngắn nhất tương ứng trong \(G\). Nếu có nhiều
  đường đi ngắn nhất thì chọn tùy ý.
  \item Tìm cây khung nhỏ nhất \(T_2\) của \(G_2\). Nếu có nhiều nghiệm thì chọn
  tùy ý.
  \item Xóa khỏi \(T_2\) các cạnh vô dụng sao cho mọi lá còn lại đều là đỉnh
  thuộc \(S\). Kết quả là cây Steiner \(T\).
\end{enumerate}

Phân tích từng bước:
\begin{itemize}
  \item Dựng \(G_1\) cần chạy đường đi ngắn nhất một nguồn \(|S|\) lần, độ phức
  tạp \(O(|S||V|^2)\) nếu dùng cách cơ bản.
  \item Tìm cây khung nhỏ nhất của \(G_1\); vì \(G_1\) là đồ thị đầy đủ, nên dùng
  Prim với độ phức tạp \(O(|S|^2)\).
  \item Khi dựng \(G_2\), mỗi đường đi ngắn nhất có thể qua \(O(|V|)\) đỉnh,
  nhưng tổng số cạnh đường đi được thêm vẫn bị chặn bởi \(O(|V|)\) theo phân
  tích của bản gốc cho trường hợp này.
  \item Tìm cây khung nhỏ nhất của \(G_2\), độ phức tạp \(O(|V|^2)\) hoặc
  \(O(|V|\log |V|)\) tùy cài đặt.
  \item Xóa cạnh vô dụng trong \(O(|V|)\).
\end{itemize}

Có thể thấy thuật toán chạy rất hiệu quả, đồng thời chất lượng xấp xỉ cũng tốt.
Phần tiếp theo chứng minh chặn trên tương đối cho nghiệm xấp xỉ của nó. Có thể
thay bước đường đi ngắn nhất bằng thuật toán hiệu quả hơn để tối ưu thêm độ
phức tạp.

\paragraph{Hiệu quả xấp xỉ.}
Gọi \(D_T\) là tổng trọng số cây Steiner \(T\) do thuật toán trên thu được, và
\(D_{\min}\) là tổng trọng số cây Steiner nhỏ nhất \(T_{\min}\). Ta sẽ chặn
tỉ số \(D_T/D_{\min}\).

\paragraph{Bổ đề 6.1.}
Với một cây \(T\) có \(m>1\) cạnh, tồn tại một vòng
\[
  u_0,u_1,\ldots,u_{2m}=u_0
\]
thỏa các tính chất:
\begin{itemize}
  \item mỗi cạnh của \(T\) xuất hiện đúng hai lần trong vòng;
  \item mỗi lá của \(T\) chỉ xuất hiện một lần trong vòng, không tính việc
  \(u_0\) lặp lại ở cuối;
  \item nếu \(u_i\) và \(u_j\), với \(i<j\), là hai lá liên tiếp trên vòng, và
  giữa chúng không có lá nào khác, thì \(u_i,u_{i+1},\ldots,u_j\) là một đường
  đi đơn trên \(T\).
\end{itemize}

\paragraph{Chứng minh.}
Chọn một đỉnh làm gốc, rồi duyệt Euler trên cây \(T\). Dãy Euler thu được thỏa
các tính chất trên nếu xem mỗi cạnh vô hướng như hai cạnh có hướng ngược nhau.

\paragraph{Định lý 6.1.}
Gọi \(\mathrm{leaf}\) là số lá của cây Steiner nhỏ nhất \(T_{\min}\). Khi đó
\[
  \frac{D_T}{D_{\min}}
  < 2\left(1-\frac1{\mathrm{leaf}}\right)
  \le 2\left(1-\frac1{|S|}\right).
\]

\paragraph{Chứng minh.}
Vì \(\mathrm{leaf}\le |S|\), bất đẳng thức bên phải là hiển nhiên.

Theo Bổ đề 6.1, trong \(T_{\min}\) có thể tạo một vòng \(L\) bằng cách đi mỗi
cạnh hai lần. Nếu trọng số vòng là \(D_L\), thì \(D_L=2D_{\min}\).

Trên vòng \(L\), chọn đường đơn dài nhất giữa hai lá liên tiếp rồi xóa nó đi;
phần còn lại là một đường đi \(P\). Vì có \(\mathrm{leaf}\) đoạn giữa các lá,
đoạn dài nhất có trọng số ít nhất \(D_L/\mathrm{leaf}\), do đó
\[
  D_P\le \left(1-\frac1{\mathrm{leaf}}\right)D_L
  =2\left(1-\frac1{\mathrm{leaf}}\right)D_{\min}.
\]
Trong \(G_1\), cạnh giữa hai đỉnh thuộc \(S\) luôn lấy đường đi ngắn nhất trong
\(G\), nên cây khung nhỏ nhất trên \(G_1\) không nặng hơn đường \(P\). Các bước
khử trùng cạnh và xóa cạnh vô dụng sau đó không làm tăng tổng trọng số. Vì vậy
\[
  D_T\le D_P
  <2\left(1-\frac1{\mathrm{leaf}}\right)D_{\min}.
\]

\paragraph{Định lý 6.2.}
Khi \(\mathrm{leaf}>2\), tỉ lệ xấu nhất đạt đúng
\[
  \frac{D_T}{D_{\min}}=
  2\left(1-\frac1{\mathrm{leaf}}\right).
\]

\paragraph{Chứng minh.}
Dựng đồ thị \(G=(V,E)\) như sau:
\[
  V=\{v_1,v_2,\ldots,v_{\mathrm{leaf}},v_{\mathrm{leaf}+1}\},
  \qquad
  S=\{v_1,v_2,\ldots,v_{\mathrm{leaf}}\}.
\]
Đặt trọng số
\[
d(v_i,v_j)=
\begin{cases}
2, & j=i+1,\ i<\mathrm{leaf},\\
1, & i\le \mathrm{leaf},\ j=\mathrm{leaf}+1,\\
\infty, & \text{các trường hợp khác.}
\end{cases}
\]

Trong trường hợp xấu nhất, thuật toán có thể chọn đường đi nối các đỉnh
\(v_1,v_2,\ldots,v_{\mathrm{leaf}}\) thành một chuỗi với tổng trọng số
\(2(\mathrm{leaf}-1)\). Trong khi đó, cây Steiner tối ưu là ngôi sao qua
\(v_{\mathrm{leaf}+1}\), có tổng trọng số \(\mathrm{leaf}\). Do đó
\[
  \frac{D_T}{D_{\min}}
  =\frac{2(\mathrm{leaf}-1)}{\mathrm{leaf}}
  =2\left(1-\frac1{\mathrm{leaf}}\right),
\]
đúng bằng chặn trên.

\subsubsection{Ví dụ 10: kiểm tra lẫn nhau đội tuyển 2015, \emph{Marketing network}}

Cho đồ thị \(G=(V,E)\). Cần tìm \(K\) rừng sinh nhỏ nhất đầu tiên, với ràng buộc
tập đỉnh đặc biệt \(S\) phải liên thông trong rừng sinh. Dữ liệu:
\[
  |V|,K<50,\quad |E|<100,\quad |S|<15,
\]
cạnh, trọng số cạnh và tập \(S\) đều được sinh ngẫu nhiên.

\paragraph{Thuật toán tối ưu.}
Bài toán nghiệm tốt thứ \(k\) có thể dùng \(A^*\). Ta mô tả trạng thái bằng việc
mỗi cạnh đã được chọn hay chưa; khi tìm kiếm, quyết định cạnh tiếp theo có được
chọn hay không.

Nếu dùng thuật toán trạng thái nén cho cây Steiner làm hàm ước lượng, ta có
\(h(s)=h^*(s)\). Mỗi lần mở rộng trạng thái tốt nhất. Vì yêu cầu là rừng sinh,
sau mỗi \(O(n)\) quyết định hiệu lực có thể tìm được nghiệm tốt nhất hiện tại;
tìm \(K\) nghiệm đầu chỉ cần \(O(Km)\) lần tính ước lượng.

Thuật toán này bảo đảm tính đúng và độ phức tạp, nhưng hiệu suất thực tế thấp,
không đủ cho miền dữ liệu của bài này.

\paragraph{Thuật toán xấp xỉ.}
Tiếp tục dùng tư tưởng \(A^*\), nhưng lấy thuật toán heuristic ở phần trước làm
ước lượng. Bất lợi là giá trị ước lượng có thể lớn hơn \(h^*(s)\), nên không còn
thể cắt tỉa một cách chắc chắn.

Ta xử lý xấp xỉ bằng cách chọn một hằng số \(c\), dùng \(c\,h(s)\) làm tham khảo
cho cận dưới, rồi cắt các trạng thái có ước lượng kém. Điều chỉnh \(c\) cho phép
cân bằng giữa tốc độ chạy và độ chính xác.

\subsection{Lời cảm ơn}

Xin cảm ơn Hiệp hội Máy tính Trung Quốc đã cung cấp nền tảng học tập và giao
lưu.

Xin cảm ơn các thầy Chen Heli, Dong Yehua, Shao Hongxiang và You Guanghui của
Trường Trung học số 1 Thiệu Hưng đã nhiều năm quan tâm và hướng dẫn tôi.

Xin cảm ơn các huấn luyện viên đội tuyển quốc gia Yu Linyun và Chen Xuji đã chỉ
dẫn.

Xin cảm ơn các anh Yu Yili, Dong Anhua và He Qizheng của Đại học Thanh Hoa đã
giúp đỡ tôi.

Xin cảm ơn các bạn Zhang Hengjie, Wang Jianhao và Jia Yuekai của Trường Trung
học số 1 Thiệu Hưng đã giúp đỡ và gợi mở cho tôi.

Xin cảm ơn các thầy cô và bạn học khác từng giúp đỡ, gợi ý cho tôi.

\subsection*{Tài liệu tham khảo của bài viết}

\begin{enumerate}[label={[\arabic*]},leftmargin=3em]
  \item Wikipedia, Heuristic (computer science).
  \item Wikipedia, Partition problem.
  \item Wikipedia, Steiner tree problem.
  \item L. Kou, G. Markowsky and L. Berman, ``A fast algorithm for Steiner trees'', \emph{Acta Informatica}.
  \item Jiang Jinye, \emph{Công cụ mạnh cho giải lặp: tối ưu và ứng dụng thuật toán SPFA}, luận văn đội tuyển quốc gia 2009.
  \item Zhou Erjin, \emph{Bàn về ứng dụng hàm ước lượng trong thi tin học}, luận văn đội tuyển quốc gia 2009.
  \item Yu Yili, \emph{Một số phương pháp tìm nghiệm tốt thứ \(k\)}, luận văn đội tuyển quốc gia 2014.
\end{enumerate}

\end{document}
